\documentclass[12pt]{article}
\usepackage[T1]{fontenc}
\usepackage[utf8]{inputenc}
\usepackage{lmodern}
\usepackage{textcomp}
\usepackage{enumitem}
\usepackage{geometry}
\usepackage{graphicx}
\usepackage{amsmath, amssymb}
\geometry{margin=1in}
\begin{document}
\begin{center}
History - Graduate Level Assessment 8 \
Total Marks: 40 \
Answer all questions.
\end{center}

\section*{Multiple Choice (10 marks)}
\begin{enumerate}[label=\arabic*.]
\item Which of the following is the Pleistocene ice mass in North America centered in the Rocky Mountains?
\begin{enumerate}[label=(\alph*)]
    \item Mississippi
    \item Rocky Mountain
    \item Appalachian
    \item McKenzie Corridor
    \item Pacifica
\end{enumerate}
\item Archaeological evidence for the collapse of civilizations suggests which of the following is the most important variable?
\begin{enumerate}[label=(\alph*)]
    \item the existence of a written language
    \item the number of years the civilization has existed
    \item the civilization's religious beliefs and practices
    \item whether warfare can be ended
    \item the size of the civilization's population
\end{enumerate}
\item The Lake Forest Archaic tradition relied on \_\_\_\_\_\_\_\_\_\_\_\_ resources; the Maritime Archaic hunted \_\_\_\_\_\_\_\_\_\_\_ creatures.
\begin{enumerate}[label=(\alph*)]
    \item pelagic; lacustrine
    \item pelagic; littoral
    \item pelagic; midden
    \item littoral; lacustrine
    \item midden; littoral
\end{enumerate}
\item A subsistence strategy and settlement pattern based on seasonality and planned acquisition of resources is known as:
\begin{enumerate}[label=(\alph*)]
    \item tactical collecting.
    \item symbolic foraging.
    \item systematic hunting.
    \item resource planning.
    \item logistical collecting.
\end{enumerate}
\item On which continent are most of the Venus figurines found?
\begin{enumerate}[label=(\alph*)]
    \item Antarctica
    \item Africa
    \item Middle East
    \item Asia
    \item Australia
\end{enumerate}
\end{enumerate}

\section*{True / False (10 marks)}
\textit{Answer True or False}
\begin{enumerate}[label=\arabic*.]
\setcounter{enumi}{5}
\item True or False: Following a long period of military setbacks against European powers, the Ottoman Empire gradually declined into the late nineteenth century. The empire allied with a different answer in the early 20 th century, with the imperial ambition of recovering its lost territories, but it dissolved in the aftermath of World War I, leading to the emergence of the new state of Turkey in the Ottoman Anatolian heartland, as well as the creation of modern Balkan and Middle Eastern states, thus ending Turkish colonial ambitions.
\item True or False: Beyoncé also ventured into acting, with a Golden Globe-nominated performance in Dreamgirls (2006), and starring roles in The Pink Panther (2006) and Obsessed (2009). Her marriage to rapper Jay Z and portrayal of Etta James in Cadillac Records (2008) influenced her third album, I Am... a different answer (2008), which saw the birth of her alter-ego Sasha Fierce and earned a record-setting six Grammy Awards in 2010, including Song of the Year for "Single Ladies (Put a Ring on It)".
\item True or False: A famous saying by the 8 th-century Indian Buddhist scholar-saint Shantideva, which a different answer often cites as his favourite verse, summarizes the Bodhisattva's intention (Bodhicitta) as follows: "For as long as space endures, and for as long as living beings remain, until then may I too abide to dispel the misery of the world.
\item True or False: They brought the Middle Chinese phonology commonly spoken in Northern China into Zhangzhou. These a different answer generally brought the language of northern Middle Chinese into the Fujian region.
\item True or False: By the 1820 s, large numbers of Jewish German women had intermarried with Christian German men and had converted to Christianity. Jewish German Eduard Lasker was a prominent German nationalist figure who promoted the unification of Germany in the mid-19 th century.
\end{enumerate}

\section*{Long Form Response (20 marks)}
\textit{Answer each question with a clear and complete explanation.}
\begin{enumerate}[label=\arabic*.]
\setcounter{enumi}{10}
\item Read the following passage and answer the question: The rebels made the first move in the war, seizing the strategic Rochester Castle, owned by Langton but left almost unguarded by the archbishop. John was well prepared for a conflict. He had stockpiled money to pay for mercenaries and ensured the support of the powerful marcher lords with their own feudal forces, such as William Marshal and Ranulf of Chester. The rebels lacked the engineering expertise or heavy equipment necessary to assault the network of royal castles that cut off the northern rebel barons from those in the south. John's strategy was to isolate the rebel barons in London, protect his own supply lines to his key source of mercenaries in Flanders, prevent the French from landing in the south-east, and then win the war through slow attrition. John put off dealing with the b... Question: Who made the first move in the war?
\item Read the following passage and answer the question: Virtually all staple foods come either directly from primary production by plants, or indirectly from animals that eat them. Plants and other photosynthetic organisms are at the base of most food chains because they use the energy from the sun and nutrients from the soil and atmosphere, converting them into a form that can be used by animals. This is what ecologists call the first trophic level. The modern forms of the major staple foods, such as maize, rice, wheat and other cereal grasses, pulses, bananas and plantains, as well as flax and cotton grown for their fibres, are the outcome of prehistoric selection over thousands of years from among wild ancestral plants with the most desirable characteristics. Botanists study how plants produce food and how to increase yields, for example thr... Question: How can the yield of food plants be increased?
\end{enumerate}

\vfill
\noindent\textit{End of Paper}
\end{document}